\begin{choice}[source={2023 秋 · 高一 · 山东烟台 · 月考校考}, topic={充分必要条件}]
下列说法错误的是
\begin{tasks}(1)
	\task ``$A \cap B = B$'' 是 ``$B = \emptyset$'' 的必要不充分条件
	\task ``$x=3$'' 的一个充分不必要条件是 ``$x^2-2x-3=0$''
	\task ``$\lvert x\rvert=1$'' 是 ``$x=1$'' 的必要不充分条件
	\task ``$m$ 是实数'' 的一个充分不必要条件是 ``$m$ 是有理数''
\end{tasks}
\end{choice}
\begin{choice-sol}
\textbf{答案 D}

本题考查充分条件与必要条件的判断.
核心方法是将命题间的逻辑关系转化为集合间的包含关系：
若 $p \Rightarrow q$，则 $p$ 对应的集合是 $q$ 对应集合的子集.

A. ``$A \cap B = B$'' 等价于 $B \subseteq A$.
``$B=\emptyset$'' 能够推出 $B \subseteq A$，但反之不成立.
因此，``$B \subseteq A$'' 是 ``$B=\emptyset$'' 的必要不充分条件.该说法正确.

B. ``$x=3$'' 对应集合 $\{3\}$.
``$x^2-2x-3=0$'' 对应解集 $\{3, -1\}$.
因为 $\{3\} \subsetneq \{3, -1\}$，
所以前者是后者的充分不必要条件.该说法正确.

C. ``$\lvert x\rvert=1$'' 对应集合 $\{1, -1\}$.
``$x=1$'' 对应集合 $\{1\}$.
因为 $\{1\} \subsetneq \{1, -1\}$，
所以前者是后者的必要不充分条件.该说法正确.

D. 命题 $p$: ``$m$ 是有理数''，对应集合 $\mathbb{Q}$.
命题 $q$: ``$m$ 是实数''，对应集合 $\mathbb{R}$.
因为 $\mathbb{Q} \subsetneq \mathbb{R}$，
所以 $p$ 是 $q$ 的充分不必要条件.
原说法颠倒了主次，故错误.
\end{choice-sol}
